\documentclass[11pt]{article}
\usepackage[margin=1in]{geometry}
\usepackage{amsmath,amssymb,booktabs}
\usepackage[colorlinks=true,linkcolor=blue,urlcolor=blue]{hyperref}

\title{Current Implemented Model: Equation Note}
\author{}
\date{\today}

\begin{document}
\maketitle

\section*{Purpose}

This note documents the model that is \emph{currently implemented in code} in
\href{/Users/fxr/Desktop/TidalLanes/data_work/src/model/spatial_equilibrium.py}{\texttt{spatial\_equilibrium.py}},
using notation aligned with the paper draft.

\section{Objects and Notation}

Let the city be partitioned into grid cells indexed by $i,j,k,l = 1,\dots,N$.

\begin{align*}
L_i &:= \text{residents living in grid } i, \\
H_j &:= \text{jobs located in grid } j, \\
M_{ij} &:= \text{commuting flow from residence } i \text{ to workplace } j, \\
t_{kl} &:= \text{direct travel time on directed edge } (k,l), \\
\tau_{ij} &:= \text{bilateral commuting cost from } i \text{ to } j, \\
n_{kl} &:= \text{lane capacity on directed edge } (k,l), \\
\phi_{kl} &:= \text{flow assigned to edge } (k,l).
\end{align*}

Total population is
\[
\bar L = \sum_i L_i.
\]

Observed objects in the current implementation are
\[
L_i^{obs}, \qquad H_j^{obs}, \qquad M_{ij}^{obs}, \qquad t_{kl}^{obs}, \qquad t_{kl}^{ff}, \qquad n_{kl}^{obs}.
\]

\section{Bilateral Commuting Costs}

The current code approximates bilateral commuting costs by shortest-path travel time:
\[
\tau_{ij} = \min_{p \in \mathcal P_{ij}} \sum_{(k,l)\in p} t_{kl},
\]
where $\mathcal P_{ij}$ is the set of feasible directed paths from $i$ to $j$.

For intrazonal commuting, the code assigns
\[
\tau_{ii}
=
\frac{1}{2}\operatorname{median}\{ t_{kl} : (k,l)\in \mathcal E,\; t_{kl}>0 \}.
\]

\section{Diagnostic Calibration}

\subsection{Diagnostic `\theta`}

The code estimates a gravity-style diagnostic elasticity from
\[
\log M_{ij}^{obs} = \mu_i + \nu_j - \theta \log \tau_{ij} + \varepsilon_{ij}.
\]

This estimate is saved as a diagnostic, but the equilibrium solver typically uses an externally supplied $\theta$.

\subsection{Diagnostic `\lambda`}

Observed OD flows are assigned to shortest paths, producing observed edge flows $\phi_{kl}^{obs}$.
Observed per-lane density is
\[
\rho_{kl}^{obs} = \frac{\phi_{kl}^{obs}}{n_{kl}^{obs}}.
\]

The diagnostic congestion regression is
\[
\log \left(\frac{t_{kl}^{obs}}{t_{kl}^{ff}}\right)
=
\eta_{\text{lane-bin}(kl)} + \lambda \log \rho_{kl}^{obs} + u_{kl}.
\]

Again, this serves as a diagnostic rather than the default structural value used in solution.

\section{Population Shares and Fundamentals}

Define resident and employment shares:
\[
\ell_i^L = \frac{L_i}{\bar L},
\qquad
\ell_j^H = \frac{H_j}{\bar L}.
\]

The code works with two fundamentals:
\[
\bar u_i^\theta
\qquad\text{and}\qquad
\bar a_j^\theta.
\]

These correspond to residential amenity and employment productivity terms in $\theta$-power form.

\section{Commuting Mass and Flows}

Given current commuting costs and current distributions, define
\[
U_i^\theta = \bar u_i^\theta (\ell_i^L)^{\beta\theta},
\qquad
A_j^\theta = \bar a_j^\theta (\ell_j^H)^{\alpha\theta}.
\]

Then commuting mass is
\[
m_{ij} = \tau_{ij}^{-\theta} U_i^\theta A_j^\theta.
\]

The corresponding OD flow is
\[
M_{ij}
=
\bar L \cdot \frac{m_{ij}}{\sum_{r,s} m_{rs}}.
\]

Hence equilibrium residents and jobs satisfy
\[
L_i = \sum_j M_{ij},
\qquad
H_j = \sum_i M_{ij}.
\]

\section{Inversion of Baseline Fundamentals}

Using observed distributions,
\[
\ell_i^{L,obs} = \frac{L_i^{obs}}{\bar L},
\qquad
\ell_j^{H,obs} = \frac{H_j^{obs}}{\bar L},
\]
the code iteratively solves
\[
\bar u_i^\theta
=
\frac{(\ell_i^{L,obs})^{1-\beta\theta}}
{\sum_j \tau_{ij}^{-\theta}\bar a_j^\theta(\ell_j^{H,obs})^{\alpha\theta}},
\]

\[
\bar a_j^\theta
=
\frac{(\ell_j^{H,obs})^{1-\alpha\theta}}
{\sum_i \tau_{ij}^{-\theta}\bar u_i^\theta(\ell_i^{L,obs})^{\beta\theta}}.
\]

\section{Route Assignment}

The current implementation uses all-or-nothing shortest-path assignment.
If $p_{ij}^*$ denotes one shortest path from $i$ to $j$, then edge flow is
\[
\phi_{kl}
=
\sum_{i,j} M_{ij}\mathbf 1\{(k,l)\in p_{ij}^*\}.
\]

\section{Congestion Equation}

This is the place where lane capacity enters the structural equilibrium.

Define observed and counterfactual per-lane density:
\[
\rho_{kl}^{obs} = \frac{\phi_{kl}^{obs}}{n_{kl}^{obs}},
\qquad
\rho_{kl}^{cf} = \frac{\phi_{kl}^{cf}}{n_{kl}^{cf}}.
\]

The code now uses free-flow edge time as the congestion primitive:
\[
t_{kl}^{cf}
=
t_{kl}^{ff}
\left(
\frac{\phi_{kl}^{cf}}{n_{kl}^{cf}}
\right)^{\lambda}.
\]

Therefore lane capacity affects equilibrium and counterfactual travel times through congestion, while $t_{kl}^{ff}$ is the infrastructure-side primitive.

\section{Nested Fixed Point}

The solver uses an outer congestion loop and an inner population--employment loop:

\begin{enumerate}
\item start from edge times $t_{kl}$,
\item compute shortest-path costs $\tau_{ij}$,
\item solve for $(L_i,H_j)$ via the commuting-flow fixed point,
\item assign edge flows $\phi_{kl}$,
\item update edge times through congestion,
\item damp the update:
\[
t_{kl}^{new,damped}
=
\delta t_{kl}^{new} + (1-\delta)t_{kl}^{old},
\]
\item iterate until convergence.
\end{enumerate}

\section{Welfare}

The scalar welfare object reported in the current code is
\[
W = \left( \sum_{i,j} m_{ij} \right)^{1/\theta}.
\]

\section{Tidal-Lane Counterfactual}

For treated directed edges, the code reallocates lane capacity as
\[
n_{kl}^{cf} = n_{kl}^{obs} + \Delta n,
\]
and, if the reverse edge exists,
\[
n_{lk}^{cf} = \max\{0.25,\; n_{lk}^{obs} - \Delta n\}.
\]

The new lane vector $n^{cf}$ is then fed back into the congestion equation and the equilibrium is resolved.

\section{Compact Summary}

The current implemented model can be summarized as
\[
\tau_{ij} = \min_{p \in \mathcal P_{ij}} \sum_{(k,l)\in p} t_{kl},
\]

\[
M_{ij} \propto \tau_{ij}^{-\theta}\bar u_i^\theta(\ell_i^L)^{\beta\theta}\bar a_j^\theta(\ell_j^H)^{\alpha\theta},
\]

\[
L_i = \sum_j M_{ij},
\qquad
H_j = \sum_i M_{ij},
\]

\[
\phi_{kl}
=
\sum_{i,j} M_{ij}\mathbf 1\{(k,l)\in p_{ij}^*\},
\]

\[
t_{kl}^{cf}
=
t_{kl}^{ff}
\left(
\frac{\phi_{kl}^{cf}}{n_{kl}^{cf}}
\right)^{\lambda}.
\]

This is the exact logic currently implemented in code.

\end{document}
